\section{Ordinary split-factor rigidity}
\label{sec:denominator}

\begin{definition}[Packet-admissible row]
\label{def:packet-admissible}
A row \(n\) is packet-admissible if the source is smooth and actual
rational compatible realizations \(\Epack_n\) and \(\Opack_n\) have been
constructed, pure of weights zero and one, with
\begin{equation}
e_n:=\rank\Epack_n=\frac{4^n+5}{3},\qquad
o_n:=\rank\Opack_n=\frac{2(4^n-4)}{3}=2(e_n-3).
\label{eq:rail-ranks}
\end{equation}
\end{definition}

The inherited construction certifies this packet data for
\(n=2,3,4\).  It does not certify semisimplicity, and it does not establish
packet-admissibility for \(n\ge5\).  Every result in this section is
conditional on Definition~\ref{def:packet-admissible}.

For a good rational prime \(p\), use geometric Frobenius and write
\[
L_p(\Vpack,u)
=\det(1-F_pu\mid\Vpack)^{-1},
\qquad
\Logz L_p(\Vpack,u)
=\sum_{m\ge1}\frac{\Tr(F_p^m\mid\Vpack)}{m}u^m.
\]

\begin{theorem}[Complete split-local classification]
\label{thm:denominator}
Let \(n\ge2\) be packet-admissible.  The following are equivalent.
\begin{enumerate}[label=\textup{(\roman*)},leftmargin=*]
\item There is an actual finite-rank rational compatible system
\(\Vpack_n\) such that, at every good rational prime \(p\) split in \(K\),
\[
\Tr(F_p\mid\Vpack_n)
=\frac4n\Tr(F_p\mid\Epack_n\oplus\Opack_n).
\]
\item There is such a system satisfying, at every good split prime,
\[
\Logz L_p(\Vpack_n,u)
=\frac4n\Logz L_p(\Epack_n\oplus\Opack_n,u).
\]
\item \(n\mid4\).
\end{enumerate}
Hence the only possible rows \(n\ge2\) are \(n=2\) and \(n=4\).
\end{theorem}

\begin{proof}
Condition (ii) implies (i) by comparing the coefficient of \(u\).
Assume (i), fix a coefficient prime \(\ell\), restrict all three
\(\ell\)-adic realizations to \(G_K\), and then take
semisimplifications.  This last operation preserves traces, characteristic
polynomials, ranks, and the purity of the two inherited rails.  In
particular, it is not an assertion that the original packet systems are
semisimple.

After enlarging to one finite coefficient extension \(E_\ell/\Ql\), choose
a common finite set \(S\) containing \(\ell\), ramification, and all bad
primes in scope.  The Grothendieck group below is generated by the
finite-dimensional continuous semisimple \(E_\ell\)-representations
obtained from these fixed-\(\ell\) realizations, unramified outside \(S\);
it is not an unrestricted representation category.

Outside finitely many excluded primes, the degree-one prime ideals of \(K\)
lie over split rational primes, and they form a Dirichlet-density-one set.
The assumed trace identity therefore holds on a dense set of Frobenius
classes for \(G_K\).  Chebotarev density
\cite[\S2.1, Theorem 1]{Serre1981} and the characteristic-zero
Brauer--Nesbitt consequence \cite{BrauerNesbitt1941} give
\begin{equation}
n[(\Res V_{n,\ell})^{\semis}]
=4[(\Res E_{n,\ell})^{\semis}]
+4[(\Res O_{n,\ell})^{\semis}]
\label{eq:k0-identity}
\end{equation}
in the semisimple Grothendieck group.

The two right-hand summands have different pure weights
\cite[\S1.2, especially (1.2.2) and (1.2.5)(i)]{Deligne1980}, so they
have no common irreducible constituent.
Separating \eqref{eq:k0-identity} by weight and then taking ranks forces
\begin{equation}
n\mid4e_n,\qquad n\mid4o_n.
\label{eq:rank-divisibility}
\end{equation}
Since \(o_n=2(e_n-3)\), these congruences imply
\(n\mid8e_n\) and \(n\mid8(e_n-3)\); subtraction gives \(n\mid24\).
The remaining finite check is
\begin{center}
\begin{tabular}{@{}c|rrrrrrr@{}}
\toprule
\(n\) & \(2\) & \(3\) & \(4\) & \(6\) & \(8\) & \(12\) & \(24\)\\
\midrule
\(4e_n\bmod n\) & \(0\) & \(2\) & \(0\) & \(2\) & \(4\) & \(8\) & \(20\)\\
\(4o_n\bmod n\) & \(0\) & \(1\) & \(0\) & \(4\) & \(0\) & \(4\) & \(16\)\\
\bottomrule
\end{tabular}
\end{center}
Only \(n=2,4\) satisfy both congruences.

Conversely, if \(n\mid4\), define the actual system
\begin{equation}
\Vpack_n=
\Epack_n^{\oplus 4/n}\oplus
\Opack_n^{\oplus 4/n}.
\label{eq:direct-copy-converse}
\end{equation}
It matches every Frobenius power trace at every good prime.  The defining
series for \(\Logz\) then proves (ii), completing the equivalence.
\end{proof}

\begin{remark}[Why total rank is insufficient]
\label{rem:total-rank}
At \(n=3\),
\[
\frac43(e_3+o_3)=\frac43\cdot63=84
\]
is integral, but the separate scaled ranks are
\[
\frac43e_3=\frac{92}{3},\qquad
\frac43o_3=\frac{160}{3}.
\]
The weight separation in \eqref{eq:k0-identity}, rather than the total
rank, is the decisive step.
\end{remark}
